\documentclass[../main.tex]{subfiles}

\begin{document}

\section{Quy trình chạy thử nghiệm hệ thống}
Quá trình chạy thử nghiệm được thực hiện theo mô hình client--server. Người dùng truy cập giao diện web tại \texttt{http://localhost:5173} (chạy local) hoặc \texttt{http://localhost:3000} (Docker), tải lên ảnh chân dung và ảnh kiểu tóc tham chiếu, sau đó gửi yêu cầu xử lý đến backend. Hệ thống tiếp nhận dữ liệu, đưa tác vụ vào hàng đợi Redis và thực hiện suy luận trên GPU trước khi trả kết quả về giao diện.

\section{Các bước vận hành chính}
Một phiên chạy thử nghiệm điển hình bao gồm:
\begin{enumerate}
    \item Người dùng tải lên ảnh chân dung đầu vào và ảnh kiểu tóc tham chiếu trên giao diện web.
    \item Frontend gọi \texttt{POST /detect-faces} để quét khuôn mặt. Nếu phát hiện nhiều khuôn mặt, hệ thống hiển thị modal FaceSelector để người dùng chọn.
    \item Frontend gửi yêu cầu \texttt{POST /generate} với ảnh đã crop, prompt và màu tóc (nếu có). Backend trả về \texttt{task\_id}.
    \item Celery worker nhận tác vụ, thực hiện tuần tự: Face Analysis $\rightarrow$ Hair \& Face Mask $\rightarrow$ Dynamic Mask Expansion $\rightarrow$ Depth Estimation $\rightarrow$ Diffusion Inference $\rightarrow$ Hair Color (tùy chọn) $\rightarrow$ Save Output.
    \item Frontend poll \texttt{GET /status/\{task\_id\}} mỗi 2 giây. Khi nhận \texttt{SUCCESS}, hiển thị ảnh kết quả.
\end{enumerate}

\section{Nội dung đánh giá}

\subsection{Kịch bản kiểm thử}

\begin{table}[H]
    \centering
    \caption{Bảng mô tả các kịch bản kiểm thử tiêu biểu}
    \label{tab:test-scenarios}
    \begin{tabularx}{\textwidth}{|>{\raggedright\arraybackslash}p{1.5cm}|>{\raggedright\arraybackslash}p{4cm}|>{\raggedright\arraybackslash}p{4cm}|X|}
        \hline
        \textbf{Mã} & \textbf{Đầu vào} & \textbf{Mục tiêu kiểm thử} & \textbf{Kết quả mong đợi} \\
        \hline
        TC01 & Khởi động đầy đủ frontend, backend, Redis và worker & Kiểm tra khả năng vận hành tổng thể & Các dịch vụ khởi động thành công, sẵn sàng tiếp nhận yêu cầu. \\
        \hline
        TC02 & Ảnh rõ nét + ảnh tóc tham chiếu & Kiểm tra quy trình tạo ảnh chuẩn & Trả về task\_id, xử lý thành công, hiển thị ảnh kết quả. \\
        \hline
        TC03 & Ảnh chân dung chính diện, nền đơn & Đánh giá bảo toàn khuôn mặt + chuyển giao tóc & Giữ đặc trưng khuôn mặt, thể hiện kiểu tóc tham chiếu. \\
        \hline
        TC04 & Ảnh góc nghiêng hoặc tóc gốc ít & Đánh giá mặt nạ động + xử lý khó & Hệ thống sinh được ảnh, chất lượng có thể giảm nhẹ. \\
        \hline
        TC05 & Ảnh kết quả + yêu cầu đổi màu & Kiểm tra phân hệ đổi màu tóc & Thay đổi màu tóc, giữ cấu trúc và độ sáng. \\
        \hline
        TC06 & Ảnh có nhiều khuôn mặt & Kiểm tra face detection + selection & Hiện modal chọn mặt, crop đúng khuôn mặt đã chọn. \\
        \hline
        TC07 & Ảnh xoay ngang hoặc lộn ngược & Kiểm tra auto-rotate & Hệ thống tự xoay ảnh $90\degree$/$180\degree$/$270\degree$, phát hiện mặt. \\
        \hline
    \end{tabularx}
\end{table}

\subsection{Đánh giá khả năng vận hành}
Kiến trúc frontend--backend--Redis--Celery hoạt động ổn định trong môi trường kiểm thử. Việc tách riêng worker xử lý AI giúp backend phản hồi nhanh ở bước tiếp nhận yêu cầu, giao diện không bị treo trong thời gian mô hình đang xử lý. Cơ chế lazy loading model giúp tiết kiệm bộ nhớ khi khởi động worker.

Quy trình auto-rotate ảnh (thử 4 góc: $0\degree$, $90\degree$, $180\degree$, $270\degree$) giúp hệ thống xử lý được cả ảnh chụp bằng điện thoại bị xoay ngang. Cơ chế face detection qua FaceSelector cho phép người dùng chọn đúng khuôn mặt trong ảnh nhóm.

\subsection{Đánh giá chất lượng kết quả}
Trên các ảnh chân dung rõ nét ở góc chính diện hoặc nghiêng vừa phải, hệ thống tạo ra kết quả trực quan, giữ được phần lớn đặc trưng khuôn mặt (hình dáng, hướng nhìn, ánh sáng) và thể hiện được kiểu tóc tham chiếu (độ dài, kiểu rẽ ngôi, độ xoăn).

Khi ảnh đầu vào có góc mặt khó, bị che khuất hoặc ánh sáng phức tạp, chất lượng mặt nạ có thể chưa chính xác hoàn toàn, dẫn đến hiện tượng lộ biên tóc hoặc sai lệch nhẹ tại vùng ghép. Hệ thống hoạt động tốt nhất với ảnh có một khuôn mặt chính, nền không quá phức tạp và vùng tóc không bị che khuất.

\section{Minh họa kết quả chạy thử}
Hình~\ref{fig:trial-result} minh họa một trường hợp chạy thử hệ thống với ảnh đầu vào và ảnh kết quả sau khi sinh kiểu tóc mới.

\begin{figure}[H]
    \centering
    \includegraphics[width=0.9\textwidth]{trial-2}
    \caption{Kết quả chạy thử nghiệm của hệ thống TryHairStyle}
    \label{fig:trial-result}
\end{figure}

\section{Bảng tổng hợp đánh giá}

\begin{table}[H]
    \centering
    \caption{Tổng hợp đánh giá hệ thống trong quá trình chạy thử nghiệm}
    \label{tab:evaluation-summary}
    \begin{tabularx}{\textwidth}{|>{\raggedright\arraybackslash}p{3.5cm}|>{\raggedright\arraybackslash}p{3cm}|X|}
        \hline
        \textbf{Tiêu chí} & \textbf{Mức đánh giá} & \textbf{Nhận xét} \\
        \hline
        Khởi động hệ thống & Đạt & Frontend, backend, Redis và worker khởi động ổn định cả local lẫn Docker. \\
        \hline
        Xử lý bất đồng bộ & Đạt & Cơ chế hàng đợi phản hồi tốt, giao diện hiển thị trạng thái xử lý. \\
        \hline
        Phát hiện khuôn mặt & Tốt & YOLO + InsightFace + AdaFace phát hiện cả khuôn mặt nghiêng, hỗ trợ auto-rotate. \\
        \hline
        Bảo toàn khuôn mặt & Khá & Đặc trưng khuôn mặt được giữ lại tốt trên ảnh rõ nét, đúng tư thế. \\
        \hline
        Chuyển giao kiểu tóc & Khá & Thể hiện được độ dài, kiểu rẽ ngôi và độ xoăn từ tóc tham chiếu. \\
        \hline
        Tính tự nhiên vùng ghép & Khá & Gaussian blur biên giảm lộ viền, nhưng một số trường hợp khó vẫn có sai lệch nhẹ. \\
        \hline
        Mặt nạ động & Khá & Mở rộng mask hợp lý cho tóc dài/phồng, xử lý được trường hợp trọc đầu. \\
        \hline
        Đổi màu tóc & Tốt & HSV transfer nhanh, 15 preset + hex tùy chỉnh, giữ kết cấu tự nhiên. \\
        \hline
        Thời gian xử lý & Trung bình & Khoảng vài giây/yêu cầu trên GPU 12GB, chưa đáp ứng thời gian thực. \\
        \hline
    \end{tabularx}
\end{table}

\end{document}
